\ifdefined\gkver\else\def\gkver{student}\fi
\documentclass[\gkver]{gaokaozhenti}
\begin{document}

\chapter{2010年全国新课标}
%% paperType: 理综物理部分

\begin{enumerate}
\item
%% number: 14
%% paperName: 2010年全国新课标第 14 题 6 分
%% typeId: 2
%% type: 多选题
%% chapter: 其他
%% point: 物理学史
%% method: 无
%% score: 6
%% degree: 650
%% duplicateId: 0
%% body:
在电磁学发展过程中，许多科学家做出了贡献。下列说法正确的是\xzanswer{AC}
\fourchoices[answer=AC]
{奥斯特发现了电流磁效应；法拉第发现了电磁感应现象}
{麦克斯韦预言了电磁波；楞次用实验证实了电磁波的存在}
{库仑发现了点电荷的相互作用规律；密立根通过油滴实验测定了元电荷的数值}
{安培发现了磁场对运动电荷的作用规律；洛仑兹发现了磁场对电流的作用规律}
%% answer: AC
%% memo:
\memoanswer{
赫兹用实验证实了电磁波的存在，B 错误。洛仑兹发现了磁场对运动电荷作用的规律，D 错误。
}

\item
%% number: 15
%% paperName: 2010年全国新课标第 15 题 6 分
%% typeId: 1
%% type: 单选题
%% chapter: 力物体的平衡
%% point: 弹力
%% method: 无
%% score: 6
%% degree: 650
%% duplicateId: 0
%% body:
一根轻质弹簧一端固定，用大小为 $F_{1}$ 的力压弹簧的另一端，平衡时长度为 $l_{1}$；改用大小为 $F_{2}$ 的力拉弹簧，平衡时长度为 $l_{2}$。弹簧的拉伸或压缩均在弹性限度内，该弹簧的劲度系数为\xzanswer{C}
\fourchoices[answer=C]
{$\frac{F_{2}-F_{1}}{l_{2}-l_{1}}$}
{$\frac{F_{2}+F_{1}}{l_{2}+l_{1}}$}
{$\frac{F_{2}+F_{1}}{l_{2}-l_{1}}$}
{$\frac{F_{2}-F_{1}}{l_{2}+l_{1}}$}
%% answer: C
%% memo:
\memoanswer{
根据胡克定律有：$F_{1}=k(l_{0}-l_{1})$，$F_{2}=k(l_{2}-l_{0})$，解得：$k=\frac{F_{2}+F_{1}}{l_{2}-l_{1}}$，C 正确。
}

\item
%% number: 16
%% paperName: 2010年全国新课标第 16 题 6 分
%% typeId: 2
%% type: 多选题
%% chapter: 功和能
%% point: 功率
%% method: 极值
%% score: 6
%% degree: 450
%% duplicateId: 0
%% body:
如图所示，在外力作用下某质点运动的 $v-t$ 图象为正弦曲线。从图中可以判断\xzanswer{AD}
\onepicture[w=5cm]{figs/16.png}
\fourchoices[answer=AD]
{在 $0\sim t_{1}$ 时间内，外力做正功}
{在 $0\sim t_{1}$ 时间内，外力的功率逐渐增大}
{在 $t_{2}$ 时刻，外力的功率最大}
{在 $t_{1}\sim t_{3}$ 时间内，外力做的总功为零}
%% answer: AD
%% memo:
\memoanswer{
根据 $P=Fv$ 和图象斜率表示加速度，加速度对应合外力，外力的功率先减小后增大，B 错误。$t_{2}$ 时刻外力的功率为零，C 错误。
}

\item
%% number: 17
%% paperName: 2010年全国新课标第 17 题 6 分
%% typeId: 1
%% type: 单选题
%% chapter: 电场
%% point: 电场线
%% method: 无
%% score: 6
%% degree: 450
%% duplicateId: 0
%% body:
静电除尘器是目前普遍采用的一种高效除尘器。某除尘器模型的收尘板是很长的条形金属板，图中直线 ab 为该收尘板的横截面。工作时收尘板带正电，其左侧的电场线分布如图所示；粉尘带负电，在电场力作用下向收尘板运动，最后落在收尘板上。若用粗黑曲线表示原来静止于 P 点的带电粉尘颗粒的运动轨迹，下列 4 幅图中可能正确的是（忽略重力和空气阻力）\xzanswer{A}
\onepicture[w=3.2cm, align=right]{figs/17.png}
\fourchoices[answer=A, ispicture=true, h=2.5cm]
{figs/17a.png}
{figs/17b.png}
{figs/17c.png}
{figs/17d.png}
%% answer: A
%% memo:
\memoanswer{
粉尘受力方向为电场线方向，从静止开始运动应该是 A 图情况，不会是 BCD 情况，A 正确。
}

\item
%% number: 18
%% paperName: 2010年全国新课标第 18 题 6 分
%% typeId: 1
%% type: 单选题
%% chapter: 力物体的平衡
%% point: 多力平衡
%% method: 无
%% score: 6
%% degree: 450
%% duplicateId: 0
%% body:
如图所示，一物块置于水平地面上。当用与水平方向成 $60^{\circ}$ 角的力 $F_{1}$ 拉物块时，物块做匀速直线运动；当改用与水平方向成 $30^{\circ}$ 角的力 $F_{2}$ 推物块时，物块仍做匀速直线运动。若 $F_{1}$ 和 $F_{2}$ 的大小相等，则物块与地面之间的动摩擦因数为\xzanswer{B}
\onepicture[w=5.5cm]{figs/18.png}
\fourchoices[answer=B]
{$\sqrt{3}-1$}
{$2-\sqrt{3}$}
{$\frac{\sqrt{3}}{2}-\frac{1}{2}$}
{$1-\frac{\sqrt{3}}{2}$}
%% answer: B
%% memo:
\memoanswer{
物体受重力 $mg$、支持力 $F_{N}$、摩擦力 $F_{f}$、已知力 $F$ 处于平衡，根据平衡条件，有：

$F\cos60^{\circ}=\mu(mg-F\sin60^{\circ})$

$F\cos30^{\circ}=\mu(mg+F\sin30^{\circ})$

联立解得：$\mu=2-\sqrt{3}$。
}

\item
%% number: 19
%% paperName: 2010年全国新课标第 19 题 3 分
%% typeId: 1
%% type: 单选题
%% chapter: 电路
%% point: 电路中的图象
%% method: 无
%% score: 3
%% degree: 450
%% duplicateId: 0
%% body:
电源的效率 $\eta$ 定义为外电路电阻消耗的功率与电源的总功率之比。在测电源电动势和内电阻的实验中得到的实验图线如图所示，图中 $U$ 为路端电压，$I$ 为干路电流，a、b 为图线上的两点，相应状态下电源的效率分别为 $\eta_{a}$、$\eta_{b}$。由图可知 $\eta_{a}$、$\eta_{b}$ 的值分别为\xzanswer{D}
\onepicture[w=4cm]{figs/19.png}
\fourchoices[answer=D]
{$\frac{3}{4}$、$\frac{1}{4}$}
{$\frac{1}{3}$、$\frac{2}{3}$}
{$\frac{1}{2}$、$\frac{1}{2}$}
{$\frac{2}{3}$、$\frac{1}{3}$}
%% answer: D
%% memo:
\memoanswer{
电源效率 $\eta=\frac{U}{E}$，$E$ 为电源的总电压（即电动势），根据图象可知 $U_{a}=\frac{2}{3}E$、$U_{b}=\frac{1}{3}E$，所以选项 D 正确。
}

\item
%% number: 20
%% paperName: 2010年全国新课标第 20 题 6 分
%% typeId: 1
%% type: 单选题
%% chapter: 曲线运动
%% point: 人造地球卫星
%% method: 无
%% score: 6
%% degree: 450
%% duplicateId: 0
%% body:
太阳系中的 8 大行星的轨道均可以近似看成圆轨道。下列 4 幅图是用来描述这些行星运动所遵从的某一规律的图像。图中坐标系的横轴是 $\lg\frac{T}{T_{0}}$，纵轴是 $\lg\frac{R}{R_{0}}$；这里 $T$ 和 $R$ 分别是行星绕太阳运行的周期和相应的圆轨道半径，$T_{0}$ 和 $R_{0}$ 分别是水星绕太阳运行的周期和相应的圆轨道半径。下列 4 幅图中正确的是\xzanswer{B}
\fourchoices[answer=B, ispicture=true, h=2.4cm]
{figs/20a.png}
{figs/20b.png}
{figs/20c.png}
{figs/20d.png}
%% answer: B
%% memo:
\memoanswer{
根据开普勒周期定律：$T^{2}=kR^{3}$，$T_{0}^{2}=kR_{0}^{3}$，两式相除后取对数，得：$\lg\frac{T^{2}}{T_{0}^{2}}=\lg\frac{R^{3}}{R_{0}^{3}}$，整理得：$2\lg\frac{T}{T_{0}}=3\lg\frac{R}{R_{0}}$，选项 B 正确。
}

\item
%% number: 21
%% paperName: 2010年全国新课标第 21 题 6 分
%% typeId: 1
%% type: 单选题
%% chapter: 电磁感应
%% point: 切割磁感线电动势
%% method: 无
%% score: 6
%% degree: 450
%% duplicateId: 0
%% body:
如图所示，两个端面半径同为 $R$ 的圆柱形铁芯同轴水平放置，相对的端面之间有一缝隙，铁芯上绕导线并与电源连接，在缝隙中形成一匀强磁场。一铜质细直棒 ab 水平置于缝隙中，且与圆柱轴线等高、垂直。让铜棒从静止开始自由下落，铜棒下落距离为 $0.2R$ 时铜棒中电动势大小为 $E_{1}$，下落距离为 $0.8R$ 时电动势大小为 $E_{2}$。忽略涡流损耗和边缘效应。关于 $E_{1}$、$E_{2}$ 的大小和铜棒离开磁场前两端的极性，下列判断正确的是\xzanswer{D}
\onepicture[w=4.2cm]{figs/21.png}
\fourchoices[answer=D]
{$E_{1}>E_{2}$，a 端为正}
{$E_{1}>E_{2}$，b 端为正}
{$E_{1}<E_{2}$，a 端为正}
{$E_{1}<E_{2}$，b 端为正}
%% answer: D
%% memo:
\memoanswer{
根据 $E=BLv$，$E_{1}=B\times2\sqrt{0.96}R\sqrt{g\times0.2R}$，$E_{2}=B\times2\sqrt{0.36}R\sqrt{g\times0.8R}$，可见 $E_{1}<E_{2}$。又根据右手定则判定电流方向从 a 到 b，在电源内部电流从电源负极流向正极，所以 D 正确。
}

\item
%% number: 22
%% paperName: 2010年全国新课标第 22 题 4 分
%% typeId: 4
%% type: 实验题
%% chapter: 功和能
%% point: DIS研究机械能守恒定律
%% method: 无
%% score: 4
%% degree: 650
%% duplicateId: 0
%% body:
图为验证机械能守恒定律的实验装置示意图。现有的器材为：带铁夹的铁架台、电磁打点计时器、纸带、带铁夹的重锤、天平。回答下列问题：
\begin{enumerate}
	\item
	为完成此实验，除了所给的器材，还需要的器材有\_\_\_\_\_\_\_。（填入正确选项前的字母）

	（A）米尺\quad（B）秒表\quad（C）$0\sim12\UV$ 的直流电源\quad（D）$0\sim12\UV$ 的交流电源
	\item
	实验中误差产生的原因有\_\_\_\_\_\_。（写出两个原因）
\end{enumerate}
\onepicture[h=5.5cm, align=right]{figs/22.png}
%% answer: 
\jdanswer{
\begin{enumerate}
	\item
	AD

	\item
	纸带与打点计时器之间有摩擦，用米尺测量纸带上点的位置时读数有误差，计算势能变化时，选取始末两点距离过近，交流电频率不稳定。

\end{enumerate}
}
%% memo:
\memoanswer{
（1）用 A 项米尺测量长度，用 D 项交流电源供打点计时器使用。

（2）纸带与打点计时器之间有摩擦，用米尺测量纸带上点的位置时读数有误差，计算势能变化时，选取始末两点距离过近，交流电频率不稳定。
}

\item
%% number: 23
%% paperName: 2010年全国新课标第 23 题 11 分
%% typeId: 4
%% type: 实验题
%% chapter: 电路
%% point: 电路实验
%% method: 无
%% score: 11
%% degree: 450
%% duplicateId: 0
%% body:
用对温度敏感的半导体材料制成的某热敏电阻 $R_{T}$，在给定温度范围内，其阻值随温度的变化是非线性的。某同学将 $R_{T}$ 和两个适当的固定电阻 $R_{1}$、$R_{2}$ 连成图 \ref{2010全国新课标23a} 虚线框内所示的电路，以使该电路的等效电阻 $R_{L}$ 的阻值随 $R_{T}$ 所处环境温度的变化近似为线性的，且具有合适的阻值范围。为了验证这个设计，他采用伏安法测量在不同温度下 $R_{L}$ 的阻值，测量电路如图 \ref{2010全国新课标23a} 所示，图中的电压表内阻很大。$R_{L}$ 的测量结果如表 1 所示。

\begin{center}\small
\begin{tblr}{|c|c|c|c|c|c|c|c|}
\hline
温度 $t$（$\UdC$） & 30.0 & 40.0 & 50.0 & 60.0 & 70.0 & 80.0 & 90.0 \\
\hline
$R_{L}$ 阻值（$\UO$） & 54.3 & 51.5 & 48.3 & 44.7 & 41.4 & 37.9 & 34.7 \\
\hline
\end{tblr}
\end{center}

\twopicture[showlabel=false, numA=1, labelA=2010全国新课标23a, numB=2, labelB=2010全国新课标23b, wA=6.5cm, wB=8cm]
{figs/23-1.png}
{figs/23-2.png}

回答下列问题：
\begin{enumerate}
	\item
	根据图 \ref{2010全国新课标23a} 所示的电路，在图 \ref{2010全国新课标23b} 所示的实物图上连线。
	\item
	为了检验 $R_{L}$ 与 $t$ 之间近似为线性关系，在坐标纸上作 $R_{L}-t$ 关系图线。
	\item
	在某一温度下，电路中的电流表、电压表的示数如图 \ref{2010全国新课标23c}、图 \ref{2010全国新课标23d} 所示。电流表的读数为\_\_\_\_，电压表的读数为\_\_\_\_\_\_。此时等效电阻 $R_{L}$ 的阻值为\_\_\_\_\_\_；热敏电阻所处环境的温度约为\_\_\_\_\_\_\_\_。
\end{enumerate}

\onepicture[w=9cm]{figs/23b.png}

\twopicture[showlabel=false, numA=3, labelA=2010全国新课标23c, numB=4, labelB=2010全国新课标23d, wA=6.5cm, wB=6.5cm]
{figs/23-3.png}
{figs/23-4.png}

%% answer: 
\jdanswer{
\begin{enumerate}
	\item
	如图所示

	\item
	如图所示

	\item
	$115.0\UmA$，$5.00\UV$，$43.5\UO$，$64.0\UdC$

\end{enumerate}
}
%% memo:
\memoanswer{
（1）根据电路图连接电路。

\onepicture[w=10cm]{figs/23_an1.png}

（2）根据数据描出点，拟合成直线。

\onepicture[w=12cm]{figs/23_an2.png}

（3）$I=\frac{23.0}{30}\times150=115.0\UmA$，$U=\frac{50.0}{60}\times6=5.00\UV$，$R=\frac{5.00}{115.0\times10^{-3}}=43.5\UO$，对照图找出相应的温度 $64.0\UdC$。
}

\item
%% number: 24
%% paperName: 2010年全国新课标第 24 题 14 分
%% typeId: 6
%% type: 计算题
%% chapter: 直线运动
%% point: 多段匀变速直线运动
%% method: 无
%% score: 14
%% degree: 450
%% duplicateId: 0
%% body:
短跑名将博尔特在北京奥运会上创造了 $100\Um$ 和 $200\Um$ 短跑项目的新世界纪录，他的成绩分别是 $9.69\Us$ 和 $19.30\Us$。假定他在 $100\Um$ 比赛时从发令到起跑的反应时间是 $0.15\Us$，起跑后做匀加速运动，达到最大速率后做匀速运动。$200\Um$ 比赛时，反应时间及起跑后加速阶段的加速度和加速时间与 $100\Um$ 比赛时相同，但由于弯道和体力等因素的影响，以后的平均速率只有跑 $100\Um$ 时最大速率的 $96\%$。求：
\begin{enumerate}
	\item
	加速所用时间和达到的最大速率；
	\item
	起跑后做匀加速运动的加速度。（结果保留两位小数）
\end{enumerate}
%% answer: 
\jdanswer{
\begin{enumerate}
	\item
	$1.29\Us$，$11.24\Ums$

	\item
	$8.71\Umsq$

\end{enumerate}
}
%% memo:
\memoanswer{
（1）设加速所用时间 $t$，匀速运动的速度为 $v$，则有：

$\frac{1}{2}vt+v(9.69-0.15-t)=100$，

$\frac{1}{2}vt+0.96v(19.30-0.15-t)=200$。

联立解得：$t=1.29\Us$，$v=11.24\Ums$。

（2）设起跑后做匀加速运动的加速度为 $a$，则 $a=\frac{v}{t}=\frac{11.24}{1.29}\Umsq=8.71\Umsq$。
}

\item
%% number: 25
%% paperName: 2010年全国新课标第 25 题 18 分
%% typeId: 6
%% type: 计算题
%% chapter: 磁场
%% point: 带电粒子在磁场中的运动
%% method: 无
%% score: 18
%% degree: 450
%% duplicateId: 0
%% body:
如图所示，在 $0\leq x\leq a$、$0\leq y\leq\frac{a}{2}$ 范围内有垂直于 $xy$ 平面向外的匀强磁场，磁感应强度大小为 $B$。坐标原点 O 处有一个粒子源，在某时刻发射大量质量为 $m$、电荷量为 $q$ 的带正电粒子，它们的速度大小相同，速度方向均在 $xy$ 平面内，与 $y$ 轴正方向的夹角分布在 $0\sim90^{\circ}$ 范围内。已知粒子在磁场中做圆周运动的半径介于 $a/2$ 到 $a$ 之间，从发射粒子到粒子全部离开磁场经历的时间恰好为粒子在磁场中做圆周运动周期的四分之一。求最后离开磁场的粒子从粒子源射出时的
\begin{enumerate}
	\item
	速度的大小；
	\item
	速度方向与 $y$ 轴正方向夹角的正弦。
\end{enumerate}
\onepicture[w=5.5cm, align=right]{figs/25.png}
%% answer: 
\jdanswer{
\begin{enumerate}
	\item
	$v=\left(2-\frac{\sqrt{6}}{2}\right)\frac{aqB}{m}$

	\item
	$\sin\alpha=\frac{6-\sqrt{6}}{10}$

\end{enumerate}
}
%% memo:
\memoanswer{
设粒子的发射速度为 $v$，粒子做圆周运动的轨道半径为 $R$，根据牛顿第二定律和洛伦兹力得：$qvB=m\frac{v^{2}}{R}$，解得：$R=\frac{mv}{qB}$。

当 $a/2<R<a$ 时，在磁场中运动的时间最长的粒子，其轨迹是圆心为 $C$ 的圆弧，圆弧与磁场的边界相切，如图所示，设该粒子在磁场中运动的时间为 $t$，依题意，$t=T/4$ 时，$\angle OCA=\pi/2$。

设最后离开磁场的粒子的发射方向与 $y$ 轴正方向的夹角为 $\alpha$，由几何关系得：

$R\sin\alpha=R-\frac{a}{2}$，$R\sin\alpha=a-R\cos\alpha$，且 $\sin^{2}\alpha+\cos^{2}\alpha=1$。

解得：$R=\left(2-\frac{\sqrt{6}}{2}\right)a$，$v=\left(2-\frac{\sqrt{6}}{2}\right)\frac{aqB}{m}$，$\sin\alpha=\frac{6-\sqrt{6}}{10}$。

\onepicture[w=7cm]{figs/25_an.png}
}

\item
%% number: 33
%% paperName: 2010年全国新课标第 33 题 5 分
%% typeId: 1
%% type: 单选题
%% chapter: 光学
%% point: 全反射
%% method: 无
%% score: 5
%% degree: 650
%% duplicateId: 0
%% body:
如图，一个三棱镜的截面为等腰直角 $\triangle ABC$，$\angle A$ 为直角。此截面所在平面内的光线沿平行于 $BC$ 边的方向射到 $AB$ 边，进入棱镜后直接射到 $AC$ 边上，并刚好能发生全反射。该棱镜材料的折射率为\_\_\_\_\_\_\_\_\_。（填入正确选项前的字母）
\onepicture[w=4.5cm, align=right]{figs/33a.png}
\fourchoices[answer=A]
{$\frac{\sqrt{6}}{2}$}
{$\sqrt{2}$}
{$\frac{3}{2}$}
{$\sqrt{3}$}
%% answer: A
%% memo:
\memoanswer{
根据折射率定义有，$\sin\angle1=n\sin\angle2$，$n\sin\angle3=1$，已知 $\angle1=45^{\circ}$，$\angle2+\angle3=90^{\circ}$，解得：$n=\frac{\sqrt{6}}{2}$。
}

\item
%% number: 33
%% paperName: 2010年全国新课标第 33 题 10 分
%% typeId: 6
%% type: 计算题
%% chapter: 机械波
%% point: 干涉衍射
%% method: 无
%% score: 10
%% degree: 450
%% duplicateId: 0
%% body:
波源 $S_{1}$ 和 $S_{2}$ 振动方向相同，频率均为 $4\UHz$，分别置于均匀介质中 $x$ 轴上的 O、A 两点处，$OA=2\Um$，如图所示。两波源产生的简谐横波沿 $x$ 轴相向传播，波速为 $4\Ums$。已知两波源振动的初始相位相同。求：
\begin{enumerate}
	\item
	简谐横波的波长；
	\item
	OA 间合振动振幅最小的点的位置。
\end{enumerate}
\onepicture[w=5.5cm, align=right]{figs/33b.png}
%% answer: 
\jdanswer{
\begin{enumerate}
	\item
	$1\Um$

	\item
	$0.25\Um$，$0.75\Um$，$1.25\Um$，$1.75\Um$

\end{enumerate}
}
%% memo:
\memoanswer{
（1）设波长为 $\lambda$，频率为 $\nu$，则 $v=\lambda\nu$，代入已知数据得：$\lambda=1\Um$。

（2）以 O 为坐标原点，设 P 为 OA 间任一点，其坐标为 $x$，则两波源到 P 点的波程差 $\Delta l=x-(2-x)$，$0\leq x\leq2$。其中 $x$、$\Delta l$ 以 $\Um$ 为单位。

合振动振幅最小的点的位置满足 $\Delta l=\left(k+\frac{1}{2}\right)\lambda$，$k$ 为整数。

解得：$x=0.25\Um$，$0.75\Um$，$1.25\Um$，$1.75\Um$。
}

\item
%% number: 34
%% paperName: 2010年全国新课标第 34 题 5 分
%% typeId: 1
%% type: 单选题
%% chapter: 光学
%% point: 光子说
%% method: 无
%% score: 5
%% degree: 650
%% duplicateId: 0
%% body:
用频率为 $\nu_{0}$ 的光照射大量处于基态的氢原子，在所发射的光谱中仅能观测到频率分别为 $\nu_{1}$、$\nu_{2}$、$\nu_{3}$ 的三条谱线，$\nu_{1}>\nu_{2}>\nu_{3}$，则\_\_\_\_\_\_\_\_\_。（填入正确选项前的字母）
\fourchoices[answer=B]
{$\nu_{0}<\nu_{1}$}
{$\nu_{3}=\nu_{1}+\nu_{2}$}
{$\nu_{0}=\nu_{1}+\nu_{2}+\nu_{3}$}
{$\frac{1}{\nu_{1}}=\frac{1}{\nu_{2}}+\frac{1}{\nu_{3}}$}
%% answer: B
%% memo:
\memoanswer{
大量氢原子跃迁时只有三个频率的光谱，这说明是从 $n=3$ 能级向低能级跃迁，根据能量守恒有，$h\nu_{3}=h\nu_{1}+h\nu_{2}$，解得：$\nu_{3}=\nu_{1}+\nu_{2}$，选项 B 正确。
}

\item
%% number: 34
%% paperName: 2010年全国新课标第 34 题 10 分
%% typeId: 6
%% type: 计算题
%% chapter: 运动和力
%% point: 动量和能量
%% method: 无
%% score: 10
%% degree: 450
%% duplicateId: 0
%% body:
如图所示，光滑的水平地面上有一木板，其左端放有一重物，右方有一竖直的墙。重物质量为木板质量的 2 倍，重物与木板间的动摩擦因数为 $\mu$。使木板与重物以共同的速度 $v_{0}$ 向右运动，某时刻木板与墙发生弹性碰撞，碰撞时间极短。求木板从第一次与墙碰撞到再次碰撞所经历的时间。设木板足够长，重物始终在木板上。重力加速度为 $g$。
\onepicture[w=6cm, align=right]{figs/34.png}
%% answer: 
\jdanswer{
$\frac{4v_{0}}{3\mu g}$
}
%% memo:
\memoanswer{
木板第一次与墙碰撞后，向左匀减速直线运动，直到静止，再反向向右匀加速直线运动直到与重物有共同速度，再往后是匀速直线运动，直到第二次撞墙。

木板第一次与墙碰撞后，重物与木板相互作用直到有共同速度，动量守恒，有：$2mv_{0}-mv_{0}=(2m+m)v$，解得：$v=\frac{v_{0}}{3}$。

木板在第一个过程中，用动量定理，有：$mv-m(-v_{0})=\mu2mgt_{1}$。

用动能定理，有：$\frac{1}{2}mv^{2}-\frac{1}{2}mv_{0}^{2}=-\mu2mgs$。

木板在第二个过程中，匀速直线运动，有：$s=vt_{2}$。

木板从第一次与墙碰撞到再次碰撞所经历的时间 $t=t_{1}+t_{2}=\frac{2v_{0}}{3\mu g}+\frac{2v_{0}}{3\mu g}=\frac{4v_{0}}{3\mu g}$。
}

\end{enumerate}

\end{document}
